\section{Selected public engagement}

\cventry{2026}{Interview, \emph{Intelligenza artificiale, Stefano Coniglio: Va governata, la formazione è fondamentale}}{\emph{BergamoNews}}{24 September 2026}{}{\publink{https://www.bergamonews.it/2026/09/24/intelligenza-artificiale-stefano-coniglio-va-governata-la-formazione-e-fondamentale/924341/}{Read online}}

\cventry{2026}{Speaker, XXII ACBGroup Convention}{``Etica della responsabilità, cultura economica e scenari geopolitici. Progettare il futuro tra sviluppo, impresa e contesto globale''}{Teatro Sociale, Bergamo, 25 September 2026}{}{Round table on the role of professionals in sustainable development and on technological transformation. With Giorgio Gori, Ferruccio de Bortoli, Francesco Occhetta, and Patrizia Bussoli; moderated by Giuseppe De Luca}

\cventry{2026}{Institutional opening, ``Archeologie letterarie e artistiche dell'intelligenza artificiale''}{}{University of Bergamo, 2 October 2026}{}{Participation as Rector's Delegate for Artificial Intelligence}

\cventry{2026}{Panelist, ``Institutional Policies and Guidelines on AI''}{Tomorrow Today in Research Italia 2026 (Elsevier)}{Università Cattolica del Sacro Cuore, Milan, 11 June 2026}{}{Panel discussion on institutional AI governance and the conditions for responsible, trustworthy, and strategically aligned AI adoption}

\cventry{2026}{Invited speaker, ``Le nuove sfide dell'intelligenza artificiale. Diritto e Responsabilità''}{SACBO / University of Bergamo}{Bergamo, 7 May 2026}{}{}

\cventry{2026}{Discussant, presentation of \emph{Algomedia} (2026)}{with Pietro Conte and Anna Caterina Dalmasso}{``Prismi -- Sei parole chiave per un'ecologia dell'intelligenza artificiale'', University of Bergamo, 7 May 2026}{}{}

\cventry{2026}{Invited speaker, ``Governare l'intelligenza artificiale. Riflessi di una rivoluzione su economia, lavoro, scuola e democrazia''}{Fondazione Gritti Minetti / Officina Democratica}{Sala Galmozzi, Bergamo, 7 May 2026}{}{With Anna Ascani, Simone Varva, and Silvia Patrizia Gadda}

\cventry{2026}{Opening talk, ONDIF Bergamo \emph{Question Time} on artificial intelligence}{ONDIF Bergamo (Osservatorio Nazionale sul Diritto di Famiglia)}{Wellness Café, Bergamo, 18 February 2026}{}{Talk for family lawyers. With Francesca Locatelli and Francesco Stabile; hosted by Simona Russo}

\cventry{2025}{Panelist, ``Innovazione e sviluppo del territorio''}{XX Congresso Territoriale CISL Bergamo}{Bergamo, 11 April 2025}{}{Contribution: ``Uno sguardo sull'IA moderna''. With Giorgio Gori, Claudia Terzi, and Maria Paola Esposito; moderated by Silvana Galizzi}

\section{Selected professional training on AI}
\cventry{2026}{Training session, ``Come l’IA ha iniziato a parlare: una sguardo storico-scientifico sull’IA moderna''}{Centro Scolastico La Traccia}{Calcinate (BG), 3 September 2026}{}{Three-hour training session for teaching staff}

\cventry{2026}{University-wide training cycle on the use of artificial intelligence}{Faculty and technical-administrative staff}{University of Bergamo}{}{Training on the use and adoption of modern AI and generative-AI tools}

\cventry{2026}{AI training seminar}{Femca CISL Bergamo}{Bergamo}{}{Training session for council members}

\cventry{2026}{Invited speaker, workshop ``Diritto di famiglia e I.A. Dalle basi alla pratica. Atti processuali e negoziazione''}{ONDIF Bergamo (Osservatorio Nazionale sul Diritto di Famiglia) and Department of Law}{University of Bergamo, 5 and 19 June 2026}{}{Two four-hour afternoon sessions, accredited with 6 training credits for lawyers. With Francesca Locatelli and Francesco Stabile; moderated by Simona Russo}

\cventry{2026}{Invited seminar on artificial intelligence}{ITEMA AI Academy}{Itema Group, Colzate (BG)}{}{}
